\section{Introduction}
Gregory Clark's \emph{To Have and Have Not} \citeyearpar{Clark2026}
completes the trilogy begun with \emph{A Farewell to Alms}
\citeyearpar{Clark2007} and \emph{The Son Also Rises}
\citeyearpar{Clark2014}. It asks two questions: ``To what extent are
social outcomes such as occupations, income, wealth, longevity, and
health socially versus genetically determined? And to what degree can
life outcomes be changed by social interventions?'' (p.~5). The book's
ambition is to answer both through the study of family resemblance
across four centuries of English history.

Its foundation is an unusual historical resource assembled jointly with
Neil Cummins: genealogies covering more than 400,000 individuals and
approximately 1.7 million church marriage records. By combining
genealogical research with census, probate, and other records, they
trace education, occupation, wealth, and other outcomes across
generations and among relatives as distant as fourth cousins
(pp.~20--26, 35--36, 53).
% SOURCE CHECK: p. 5 reports 449,000 individuals; p. 20 reports 436,185.
% Rounded counts avoid choosing between the manuscript's inconsistent totals.

Clark advances two main empirical claims. First, after accounting for
measurement error, social status is remarkably persistent across
generations, with an underlying intergenerational correlation near 0.8
that has barely changed over four centuries. Second, the correlations
in social outcomes across different types of relatives---from
parent--child pairs and siblings to fourth cousins---``fit
extraordinarily well'' a simple model of additive genetic inheritance
with strong assortative mating (pp.~6, 53).

Clark recognizes that cultural transmission could, in principle,
produce the same correlations (p.~88). The second part therefore seeks
evidence that distinguishes the competing explanations. Comparisons
involving family size, birth order, birthplace, parental death, and
other circumstances test what he calls the genetic model's ``clear
predictions'' (p.~7). He interprets their results as evidence that
direct genetic transmission predominates and social influences
contribute little to the persistence of status (pp.~8--9).

The third part develops the historical and policy implications Clark
draws from this interpretation. He argues that ``modern societies spend
too much on education'' (p.~9), favouring income transfers and wage
compression as more effective ways to improve the circumstances of the
disadvantaged (p.~217). He also assigns migration, differential
fertility, and marital assortment a central role in shaping the
distribution of abilities and societies' long-run prospects
(pp.~217--218).

Clark deserves considerable credit for this imaginative data-collection
effort and the interesting empirical regularities it documents. I am
far less persuaded by his interpretation and conclusions. The estimates
rest on strong exclusions of shared environmental influences and on
mating and measurement assumptions that a close fit cannot validate.
This identification problem recalls the IQ debates of the 1970s
\citep{Loehlin1978,Goldberger1978}. Applying continuous-trait formulas to
binary outcomes introduces a further source of misleading estimates,
while Part~II's tests do not consistently distinguish genetic from
social transmission.

Modern genomics offers independent evidence on these assumptions.
Random genetic transmission within families helps separate genetic
effects from family circumstances; existing findings already challenge
treating population genetic associations as direct causal effects
\citep{YoungEtAl2019,YoungEtAl2022}. These methods offer a way to test
explanations that kinship correlations alone cannot distinguish.

Even accepting the estimates leaves the objections of
\citet{Goldberger1979} and \citet{Jencks1980}. Goldberger's point concerns
what is being estimated: heritability is an $R^2$ describing observed
variation, not the response to an intervention. Jencks's concerns the
mechanism: genetic effects can operate through social responses to
inherited characteristics. Nor does unchanged relative mobility imply
that policies failed to improve living standards. The review proceeds
from the model and its estimation to the environmental and genomic
evidence, then examines the historical and policy conclusions in light
of these distinctions.
